\documentclass[11pt,letterpaper]{article}
\usepackage{hanzo-defense}

\doctitle{Hanzo AI for Compliance:\\An OSS Cloud Built to Pass the Audit}
\docid{HAI-2026-016}
\docversion{Version 1.0}
\docdate{2026}
\docauthors{Zach Kelling --- Hanzo Team}
\docorgs{Hanzo Industries\\research@hanzo.ai}

\begin{document}
\makecover

\begin{abstract}
Compliance is, in operational terms, an evidence-production problem. An auditor
for SOC~2 or ISO~27001, an examiner for FINRA/SEC, a HIPAA assessor, or---for a
government system---an assessor for an Authority to
Operate (ATO) does not primarily ask whether a system \emph{could} be secure; they
ask the operator to \emph{produce the artifacts} that demonstrate a control is
designed and operating---access reviews, key-rotation records, immutable logs,
change tickets, configuration baselines. Most stacks bolt auditing on after the
fact, so producing those artifacts becomes months of manual screenshot-gathering,
log-stitching, and spreadsheet reconciliation repeated at every cycle. This paper
argues that Hanzo's sovereign open-source stack inverts that cost structure by
emitting the evidence auditors ask for \emph{by construction}: a centralized
identity and access layer (IAM) produces access-control evidence natively; a key
management service (KMS) eliminates plaintext secrets and produces key-custody
evidence; a cryptographically signed, on-chain command/audit log produces
tamper-evident, non-repudiable audit records; reproducible builds and policy-as-code
produce change-management evidence; and 50 machine-checked Lean~4 proofs plus NIST
Known-Answer-Test validation of \fips{203/204/205} produce machine-checked
cryptographic assurance rather than ``we tested it.'' We give the control-framework
mapping from these native capabilities to the SOC~2 Trust Services Criteria, to
FINRA/SEC recordkeeping, and to the NIST SP~800-53 control families (AC, AU, IA,
SC, CM) that underpin RMF and the ATO process. We then report an anonymized case
study in which this same stack carried a regulated securities platform through
FINRA/SEC approval and SOC~2 with a small team---while a $200{+}$-microservice to
3-binary consolidation shrank the very surface that had to be audited. Throughout,
claims are kept careful: the stack \emph{provides evidence for} and \emph{maps to}
controls; certification is an assessor's judgment, not a vendor's.
\end{abstract}

\paragraph{Executive Summary.}
Passing an audit is mostly about producing proof. The painful part of SOC~2,
FINRA/SEC, or a defense ATO is not deciding to be secure---it is assembling the
mountain of evidence that shows, control by control, that the system is secure and
stays that way: who can access what, where the secrets live and how often they
rotate, an unalterable record of every privileged action, proof that the code in
production is the code that was reviewed. When auditing is an afterthought, that
evidence has to be reconstructed by hand every cycle, which is why compliance is
expensive and recurring. Hanzo's stack is built so the evidence is a byproduct of
running the system. Identity is centralized, so access reviews and least-privilege
are queries, not archaeology. Secrets live in a key vault and are never in
plaintext, so key-custody evidence exists by default. Every privileged action is
cryptographically signed and written to an immutable, tamper-evident log, so the
audit trail cannot be quietly edited---and the same property satisfies
financial-recordkeeping rules that demand write-once records and supports
non-repudiation. Builds are reproducible and policy is code, so change-control
evidence is generated automatically. And the cryptography is not merely tested but
\emph{proved}---machine-checked theorems re-verified on every commit. The net effect
is that SOC~2, ISO~27001, FINRA/SEC, HIPAA, and---for government systems---the
RMF/ATO process all become tractable: less code to
assess, a smaller audit boundary, and an evidence pipeline that is collected
continuously rather than re-earned from scratch each year. For any regulated
enterprise or government program---defense included---modernizing onto this stack
means inheriting the audit pipeline, not building it.

\tableofcontents
\newpage

% ============================================================================
\section{Compliance Is an Evidence-Production Problem}
% ============================================================================

A control framework---SOC~2, FINRA/SEC recordkeeping, NIST SP~800-53 under the Risk
Management Framework (RMF)---is a catalog of assertions a system is required to
satisfy. An audit or assessment is the exercise of \emph{substantiating those
assertions with artifacts}. The decisive distinction in practice is not between
secure and insecure systems; it is between systems that can \emph{produce the
evidence} of their controls on demand and systems that cannot. An assessor's day is
spent asking ``show me'': show me the list of who has administrative access, show me
that it was reviewed last quarter, show me that this secret was rotated, show me an
unalterable record that this privileged action happened and who performed it, show
me that the binary running in production is the artifact your pipeline built from
reviewed source.

When a stack is assembled from many independently-operated services with logging and
identity bolted on afterward, answering ``show me'' is manual archaeology. Access
lists are scattered across dozens of consoles; secrets are pasted into environment
variables and config files; logs are mutable application logs spread across services
with no integrity guarantee; the relationship between source, review, and the
deployed artifact is reconstructed from memory. The result is the familiar
compliance cost: a SOC~2 Type~II observation window plus weeks of evidence
collection; a FINRA/SEC examination that demands recordkeeping the architecture was
never designed to produce; an ATO package whose body of evidence (the security
controls assessment) is assembled by hand. And because nothing about the system
\emph{emits} this evidence, the cost recurs in full at every audit cycle and every
re-authorization.

The thesis of this paper is structural: if the stack is built so that the evidence
auditors ask for is a \emph{native output} of normal operation, the cost collapses.
The audit stops being a discovery project and becomes a query against artifacts the
system already produces. Hanzo's stack is built this way deliberately, and the rest
of this paper walks the specific capabilities, maps them to the frameworks, and
shows---with an anonymized production case---what it does to the audit.

% ============================================================================
\section{Evidence by Construction}
% ============================================================================

The stack is consolidated around a small number of owned, open-source components.
Each one, as a side effect of doing its job, emits a class of evidence that an
auditor asks for. We take them in turn and name the control objective each supports.

\subsection{Identity and Access (IAM): access-control evidence}

All authentication and authorization route through one identity layer (Hanzo IAM)
rather than per-service credentials. Authentication is centralized
(OIDC/OAuth~2.1 with PKCE; one canonical set of endpoints); authorization is
role- and scope-based; and because identity is a single source of truth rather than
a federation of ad-hoc accounts, least-privilege is enforceable and \emph{auditable}.
The evidence an assessor wants for access control---the complete roster of
principals and their entitlements, the record of access grants and revocations, the
periodic access review---are queries against one system rather than a reconciliation
across many.

This supports the SOC~2 Security (Common Criteria) logical-access criteria and maps
directly to the NIST SP~800-53 \textbf{AC} (Access Control) and \textbf{IA}
(Identification and Authentication) families: centralized account management, enforced
least privilege, separation of duties expressed as roles, and unique attributable
identities for every actor. Because there is one identity plane, an access review is
\emph{produced}, not assembled.

\subsection{Secrets and Keys (KMS): confidentiality and key-management evidence}

Secrets never live in plaintext. A key management service (Hanzo KMS) holds
credentials and cryptographic keys under centralized custody with rotation, and
services obtain secrets at runtime rather than embedding them. There are no
long-lived secrets in source, in container images, or in configuration. The evidence
auditors request for confidentiality and key management---where keys live, who can
reach them, how custody is separated, when each was rotated---is the KMS's normal
bookkeeping.

This supports the SOC~2 Confidentiality criteria and maps to the SP~800-53
\textbf{SC} (System and Communications Protection) family, specifically the
key-establishment and key-management controls (SC-12/SC-13) and protection of
secrets at rest. The structural property---\emph{no plaintext secret anywhere in the
pipeline}---is itself the control, and the KMS's rotation log is the evidence that it
operates. A verification gate enforces the property mechanically: the build fails if a
plaintext-secret pattern reappears, so ``no secrets in code'' is checked rather than
asserted.

\subsection{Immutable audit trail: audit-logging, integrity, and non-repudiation}

This is the keystone. Every privileged action is cryptographically signed by the
identity that performed it and written to an \emph{append-only, on-chain,
tamper-evident} command/audit log. Two properties follow that ordinary application
logging cannot provide:

\begin{itemize}
  \item \textbf{Tamper-evidence.} The log is hash-chained and consensus-anchored, so
  any retroactive edit, deletion, or reordering is detectable. An auditor does not
  have to trust that the log was not altered; the log's own structure demonstrates it.
  \item \textbf{Non-repudiation.} Because each entry is signed by the acting
  identity (post-quantum signatures; \S\ref{sec:pq}), the actor cannot later deny the
  action, and a reviewer cannot forge one. The signature binds \emph{who} to
  \emph{what} and \emph{when}.
\end{itemize}

This maps to the SP~800-53 \textbf{AU} (Audit and Accountability) family in full---
audit event generation (AU-2), content of records including actor and outcome (AU-3),
protection of audit information from modification (AU-9), and non-repudiation (AU-10)---
and to the SOC~2 Security monitoring criteria and Processing Integrity. Crucially,
the same immutability is exactly what financial recordkeeping rules require: the
write-once, non-rewriteable, non-erasable (``WORM-like'') retention demanded for
broker-dealer books and records is an emergent property of an append-only,
hash-chained ledger, not a bolt-on archival appliance. One mechanism satisfies the
defense integrity control \emph{and} the securities recordkeeping rule.

\subsection{Post-quantum identity and transport: forward-looking confidentiality}
\label{sec:pq}

Identities, signatures on the audit log, and transport between services are
post-quantum (ML-KEM for key establishment, ML-DSA for signatures; \fips{203/204}).
The compliance-relevant point is durability of evidence: an audit record, a signed
action, or an encrypted channel captured today must remain confidential and
verifiable for the full retention horizon---which, for financial records and for
classified material, is measured in years to decades. ``Harvest now, decrypt later''
makes classical-only confidentiality a liability for any record with a long required
life. Post-quantum identity and transport make the audit trail's signatures and the
data's confidentiality resistant to a future quantum adversary, which is the
forward-looking posture that CNSA~2.0 mandates for national-security systems and that
a prudent long-retention financial archive should adopt. This supports the SC family's
transmission-confidentiality and cryptographic-protection controls with algorithms
chosen for the record's lifetime, not just today's threat.

\subsection{Formal verification: machine-checked evidence, not ``we tested it''}

Most software is trusted because it was tested and nothing broke. The cryptographic
core here goes further: its security properties are proved as theorems and re-checked
by a machine on every code change. The suite comprises \textbf{50 machine-checked
Lean~4 proofs}, including \textbf{ML-DSA EU-CMA} unforgeability (\fips{204}) and
\textbf{ML-KEM IND-CCA2} security (\fips{203}), alongside consensus safety/liveness,
threshold-signature, and protocol-correctness proofs; complementary
\textbf{NIST Known-Answer-Test (KAT)} validation covers all three post-quantum
standards (\fips{203/204/205}). For an assessor this is a different \emph{kind} of
evidence: not a test report asserting that some inputs produced expected outputs, but
a machine-checked proof with an explicit, inspectable set of assumptions. A proof
regression blocks the build, so the assurance is continuous rather than a point-in-time
snapshot. This is the strongest available evidence for the SC family's cryptographic
controls and for the ``correctly implemented'' burden in any cryptographic
accreditation, and it is auditable end-to-end precisely because the stack is owned and
open---there is no third-party black box whose internals are out of the prover's reach.

\subsection{Reproducible builds, policy-as-code, and change control}

The path from source to running artifact is itself evidence. Builds are reproducible
and produced by a shared CI/CD pipeline; images are pinned to immutable,
content-addressed tags (a semver tag or a commit-pinned \texttt{sha-} digest), never
to a floating \texttt{:latest}; and deployment policy and infrastructure are
expressed as code under version control. The consequence is that the relationship
auditors probe for change management---reviewed source $\rightarrow$ built artifact
$\rightarrow$ deployed instance---is a chain of immutable, attributable records rather
than an after-the-fact reconstruction. Who changed what, who reviewed it, and exactly
which artifact reached production are all recorded by the normal pipeline.

This maps to the SP~800-53 \textbf{CM} (Configuration Management) family---baseline
configuration (CM-2), change control (CM-3), and configuration of the deployed
system---and to the SOC~2 change-management criteria. Because policy is code, the
configuration baseline is the repository state, and any drift is a diff.

% ============================================================================
\section{Control-Framework Mapping}
% ============================================================================

Table~\ref{tab:mapping} maps each native capability to the SOC~2 Trust Services
Criteria it provides evidence for, the NIST SP~800-53 control families it supports
under RMF/ATO, and the corresponding FINRA/SEC recordkeeping relevance. The verbs are
deliberate: a capability \emph{provides evidence for} or \emph{maps to} a control.
Whether a control is satisfied is an assessor's determination on the specific
deployment; the stack's contribution is to make the evidence native and the boundary
small.

\begin{table}[H]
\centering
\small
\begin{tabular}{@{}p{3.0cm}p{3.1cm}p{2.0cm}p{4.0cm}@{}}
\toprule
\textbf{Native capability} & \textbf{SOC~2 TSC} & \textbf{800-53 family} & \textbf{FINRA/SEC \& ATO relevance} \\
\midrule
Centralized IAM (OIDC; RBAC; least privilege; access reviews)
& Security (CC6 logical access)
& AC, IA
& Supervisory access controls; attributable identities for the ATO access-control assessment \\
\addlinespace
KMS (no plaintext secrets; central custody; rotation)
& Confidentiality
& SC (SC-12/13)
& Protection of customer/non-public data; key-management evidence in the SCA \\
\addlinespace
Immutable, signed, on-chain audit log
& Security (monitoring); Processing Integrity
& AU (AU-2/3/9/10)
& WORM-like books-and-records retention; non-repudiation of privileged actions \\
\addlinespace
Post-quantum identity \& transport (ML-KEM / ML-DSA)
& Confidentiality
& SC (SC-8/12/13)
& Long-retention record durability; CNSA~2.0 forward posture \\
\addlinespace
Formal verification (50 Lean~4 proofs; NIST KAT)
& Security; Processing Integrity
& SC, SA
& Machine-checked ``correctly implemented'' evidence for the SCA \\
\addlinespace
Reproducible builds; policy-as-code; pinned artifacts
& Security; Availability (change mgmt)
& CM (CM-2/3)
& Change-control records; configuration baseline as code \\
\addlinespace
Consolidation (fewer services / smaller boundary)
& All five (reduced scope)
& All families (smaller boundary)
& Smaller examination surface; smaller authorization boundary \\
\bottomrule
\end{tabular}
\caption{Native stack capabilities mapped to SOC~2 Trust Services Criteria, NIST
SP~800-53 control families (RMF/ATO), and FINRA/SEC recordkeeping relevance. ``Maps
to / provides evidence for''---not a certification claim.}
\label{tab:mapping}
\end{table}

\paragraph{Five criteria, one stack.} The SOC~2 Trust Services Criteria are Security,
Availability, Confidentiality, Processing Integrity, and Privacy. The capabilities
above touch all five: \emph{Security} via IAM, the audit trail, and reproducible
builds; \emph{Confidentiality} via KMS and post-quantum transport; \emph{Processing
Integrity} via the tamper-evident log and machine-checked proofs of the logic that
moves records; \emph{Availability} via policy-as-code change management and a
consolidated, self-hostable footprint; and \emph{Privacy} via centralized
access control and confidentiality of non-public data. Because one set of mechanisms
underlies all five criteria, the evidence collected once is reused across the report
rather than re-gathered per criterion.

\paragraph{The same evidence translates.} The reason one stack serves SOC~2,
ISO~27001, FINRA/SEC, and HIPAA---and, for government systems, an ATO---
simultaneously is that the underlying control objectives
overlap heavily---access control, confidentiality, audit integrity, change control,
cryptographic correctness---and differ mainly in the framework's vocabulary and the
assessor's authority. An append-only signed ledger is ``AU-9 protection of audit
information'' to an RMF assessor and ``non-rewriteable record retention'' to a
securities examiner; it is the same ledger. Centralized least-privilege identity is
``AC-2/AC-6'' to one and ``supervisory and access controls'' to another; it is the
same IAM. Producing the evidence once and presenting it through each framework's lens
is the efficiency the consolidated stack creates.

% ============================================================================
\section{Case Study: One Stack Through FINRA/SEC and SOC~2}
% ============================================================================

The argument is not theoretical. The same stack carried a regulated securities
platform---an SEC-registered alternative trading system (ATS) with broker-dealer and
transfer-agent functions---through FINRA/SEC approval and SOC~2, with a small team, in
record time. Two facts from that engagement matter for the thesis of this paper.

\paragraph{The evidence was native.} The regulated entity ran on this stack's
centralized IAM, its KMS (no plaintext secrets), and its immutable signed audit log.
The recordkeeping obligations that fall on a broker-dealer and transfer agent---
non-rewriteable books and records, an attributable trail of who did what to which
record, controlled and reviewed access---were satisfied by capabilities the platform
already had, not by a separate compliance system layered on top. The audit and the
regulatory approval drew on artifacts the running system produced as a matter of
course.

\paragraph{Consolidation shrank what had to be audited.} In parallel, the platform's
architecture was consolidated from \textbf{$200{+}$ microservices into 3 compiled Go
binaries} (plus a small set of OSS platform services), collapsing roughly
\textbf{3.07~million lines of code to about 1.06~million}---\emph{$\sim$3.5~million
lines of dead code deleted}, around a 65--70\% reduction. The consolidated rebuild
also remediated the security-finding backlog of the legacy estate (on the order of
90+ findings) at the architectural level. This is the second, quieter half of why the
audit was tractable: \emph{every line of code and every service is attack surface that
an auditor or accreditor must reason about}. Cutting the codebase by two-thirds and
the runtime from a sprawl of services to three binaries shrinks the examination
surface, the SOC~2 system description, and---for a government system---the
authorization boundary. Less code plus centralized IAM/KMS/audit is a smaller,
cheaper, faster audit. The consolidation and the evidence-by-construction property are
the same lever viewed from two sides: one reduces \emph{what} must be proven, the other
makes the proof \emph{native}.

\medskip
\noindent\emph{Anonymization.} The regulated entity is intentionally unnamed; the
point is the stack's properties, not the client. What transfers to any future program
is the mechanism: a small team reached a regulated, audited production posture quickly
because the stack produced the evidence and the consolidation shrank the surface.

% ============================================================================
\section{Continuous Compliance}
% ============================================================================

A SOC~2 Type~II report covers an observation window; an ATO is not a one-time event
but a posture that must be \emph{maintained} under continuous monitoring; a
broker-dealer's recordkeeping is examined repeatedly. The expensive failure mode is
treating each cycle as a fresh discovery project. Because this stack emits evidence
continuously, the posture is maintained rather than re-earned:

\begin{itemize}
  \item \textbf{Always-on collection.} The audit log is generated by operation, not
  by an end-of-period scramble; access reviews are queries against live IAM state;
  key-rotation records accrue from the KMS; build and deployment records accrue from
  the pipeline. The body of evidence is a continuous stream.
  \item \textbf{Continuous verification.} The Lean proofs and KAT vectors run on every
  commit, so cryptographic assurance is re-established automatically and a regression
  blocks deployment. Policy-as-code means configuration drift surfaces as a diff, not
  as an audit finding months later.
  \item \textbf{Agentic evidence assembly.} The same agentic engineering pipeline that
  built the consolidated system can collect, organize, and map evidence to controls
  continuously---turning the ongoing-authorization (cATO) ideal into routine output
  instead of a periodic mobilization.
\end{itemize}

The effect on the ATO process specifically is that re-authorization draws on an
evidence pipeline that never stopped running. The authorization boundary is smaller
(consolidation), the controls produce their own evidence (by construction), and the
evidence is current (continuous), which is precisely the posture continuous
authorization is meant to reward.

% ============================================================================
\section{Why It Matters Across Regulated Industries and Government}
% ============================================================================

Any organization modernizing onto this stack---a healthcare provider, a financial
firm, a critical-infrastructure operator, or a government contractor or agency
(defense included)---inherits the evidence pipeline
rather than building one. The same properties that make a commercial SOC~2,
ISO~27001, FINRA/SEC, or HIPAA audit tractable are the properties a government
RMF assessment and an ATO
package require: centralized least-privilege identity (AC/IA), no plaintext secrets
and managed keys (SC), a tamper-evident non-repudiable audit trail (AU), reproducible
builds and configuration-as-code (CM), and machine-checked cryptographic correctness
(SC/SA). Because the stack is sovereign, open-source, and self-hostable, it can be
deployed in the environments these regimes demand---on-premises, air-gapped, at the
edge---without renting from a closed vendor, and an assessor can inspect the entire
path the evidence takes rather than trusting a black box. The cryptography is
post-quantum, so the records and identities survive the retention horizon that
financial, healthcare, and government archives all require.

This compliance posture is one facet of a broader sovereign stack built by Hanzo
AI---an American-owned applied-AI lab and venture studio (Techstars~'17, founded
2014). Hanzo owns both halves of the modern trust problem: the frontier AI models
(the \emph{Zen} family) and the production post-quantum cryptography that secures and
signs the systems they run in. This matters for who a regulated or accredited operator
can build on. U.S.\ frontier AI has consolidated into a closed commercial duopoly,
while the broadly \emph{open} frontier model weights available today are largely of
foreign origin---an untenable basis for a system that must be inspected, accredited,
and run on sovereign infrastructure. Hanzo is the American open-source alternative
that owns the model and the cryptography together, deployable without renting from the
closed duopoly or depending on foreign open weights. The evidence-by-construction
property documented here is the auditable backing for that claim: not a vendor
assurance an accreditor must take on faith, but artifacts---signed logs, key-custody
records, reproducible builds, machine-checked proofs---an accreditor can read.

% ============================================================================
\section{Honest Scope}
% ============================================================================

This is a capability paper, and the credible version is explicit about what the stack
does versus what an assessor does. No overstatement:

\begin{itemize}
  \item \textbf{The stack provides evidence; it does not self-certify.} SOC~2 attestation
  is issued by a CPA firm, a FINRA/SEC approval is a regulator's decision, and an ATO is
  granted by an Authorizing Official. The stack's contribution is to make the evidence
  native, current, and inspectable, and to shrink the boundary being judged---which is
  what makes those determinations faster and cheaper, not automatic.
  \item \textbf{Mapping is not equivalence.} A capability \emph{provides evidence for}
  or \emph{supports} a control family; closing a specific control still requires the
  operator to configure, document, and operate it for the assessed deployment. The
  tables here are a starting correspondence, not a completed controls matrix.
  \item \textbf{Formal verification has a stated boundary.} The 50 Lean proofs rest on
  explicit cryptographic hardness axioms (standard practice), and model checking and
  protocol analysis are on a published roadmap rather than complete. The assurance
  boundary is the true one, not a marketing one.
\end{itemize}

The strategic point holds regardless of where any one line sits: because Hanzo owns
and open-sources every layer the evidence runs through, the artifacts an auditor or
accreditor needs are produced by the system itself and can be inspected end-to-end---
which is the difference between an audit that is a discovery project and one that is a
query.

% ============================================================================
\section{Conclusion}
% ============================================================================

Compliance is an evidence-production problem, and most stacks make it expensive by
producing no evidence until an auditor asks. Hanzo's sovereign open-source stack
inverts that: centralized IAM emits access-control evidence, KMS emits key-custody
evidence, an immutable signed on-chain log emits tamper-evident and non-repudiable
audit evidence, reproducible builds and policy-as-code emit change-control evidence,
and 50 machine-checked proofs plus NIST KAT validation emit cryptographic assurance
that is proved rather than asserted. The same evidence maps across SOC~2's five Trust
Services Criteria, ISO~27001, FINRA/SEC recordkeeping, HIPAA, and the NIST~800-53
families that
underpin RMF and the ATO process---so one stack serves the commercial, financial,
healthcare, and government frameworks (defense included) at once. A production case
shows it carrying a regulated securities
platform through FINRA/SEC and SOC~2 with a small team, while a $200{+}$-microservice
to 3-binary consolidation shrank the surface that had to be audited (and, for a
government system, accredited).
For any regulated enterprise or government program---defense included---the
consequence is direct: modernize on this
stack and inherit the audit pipeline---sovereign, auditable, post-quantum, and
continuous---rather than rebuild it from scratch each cycle.

\end{document}
